% page 72 of the facsimile (image p-071 of the scan)
% block 154.9 x 254.9 mm ; left margin 8.0 mm
\pgc{-12.700mm}{0.000mm}{
\l{\cel{6}64}
\l{}
\l{\sou{bibliofilo}. Persono qua prizas e kolektas la libri editita}
\l{\cel{3}pri ta o ca temo, la diversa edituri, lia valoro, lia}
\l{\cel{3}skarseso-grado, e c. - DEFIRS.}
\l{}
\l{\sou{bibliomanio}. Prizego pasionoza di la libri. - DEFIRS.}
\l{}
\l{\sou{biblioteko}. I. Moblo garnisita ye tabuli sur qui onu ordinas}
\l{\cel{3}libri. - II. Edifico qua kontenas kolektajo de libri.}
\l{\cel{3}- DEFIRS.}
\l{}
\l{\sou{bicepso}. (\sou{anat}.) Muskulo qua havas du ligo-punti ye lua}
\l{\cel{3}parto supra (kom ex.: la muskulo avana di la avanbrakio.}
\l{\cel{3}- DEFIS.}
\l{}
\l{\sou{biciklo}. Velocipedo kun du roti inter-egala, dispozita}
\l{\cel{3}ica dop ita. - DEFIRS.}
\l{}
\l{\sou{bideto}.\cel{2}Tualeto-moblo,kun kuveto sidilatra, por ablucioni}
\l{\cel{3}intergamba (por la mulieri). - DEFIRS.}
\l{}
\l{\sou{bidono}. Lada vazo klozita, por oleo, petrolo, benzino, e c.}
\l{\cel{3}qua kontenas de l til l0 litri ; gurdatra kontenilo ek}
\l{\cel{3}lado, quan la soldato portas. - FIRS.}
\l{}
\l{\sou{bielo}. (\sou{tekn.)}\cel{2}Stango qua ligas un peco ad altra (en mashi-}
\l{\cel{3}no), generale stango di pistono a manivelo, e transformas}
\l{\cel{3}movo rektalinea a movo rotaca. - DEFIRS.}
\l{}
\l{\sou{bifsteko.}\cel{2}Loncho de bovo-karno destinita a koqueso sur}
\l{\cel{3}grililo. - DEFIRS.}
\l{}
\l{\sou{bifurkar}. I. (\sou{trans.}) Dividar a du branchi. - II. (\sou{netrans.})}
\l{\cel{3}Dividar su quale la gambi di homo bifurkas de trunko.}
\l{\cel{3}- DEFIS.}
\l{}
\l{\sou{bigama}. I. (\sou{Yuro ekleziala}).Mariajita dufoye ; mariajita kun}
\l{\cel{3}vidvino. - II. Qua su mariajis duesmafoye, la maniajo}
\l{\cel{3}unesma ne esante dissolvita. - DEFIRS.}
\l{}
\l{\sou{bigoto}.\cel{2}Qua agas devocaji streta-mente, minucioze e kelke}
\l{\cel{3}hipokrite. - DEFI.}
\l{}
\l{\sou{bikursalo}. (\sou{geom.)}\cel{3}Kurvu sur qua selektesis punto de ube}
\l{\cel{3}duktesas rekto ; ta kurvo renkontresas da ta rekto en}
\l{\cel{3}nur du punti ; la duopla punto esinte selektita origine.}
\l{\cel{3}(Ex.: lemniskato). - DEFl.}
\l{}
\l{\sou{bikuspida}. (\sou{Natur-historio}). Qua finas per du pinti. - DEFIS.}
\l{}
\l{\sou{bilo}. Liquido bitra, flava verdatre, sekrecita da hepato,,}
\l{\cel{3}qua penetras en la duodenumo segun ke la nutro-materio}
\l{\cel{3}decensas aden lu, e pri qua onu konjektas ke lu helpas}
\l{\cel{4}digestado. - DEFlS.}
\l{}
\l{\sou{bibliografio}. Enumero di la libri publikigita pri ta o ca te-}
\l{\cel{3}mo, lia imprimo-dati diversa, lia valoro, lia skarseso-}
\l{\cel{3}o \textquotedbl{}\cel{4}\textquotedbl{}\cel{2}+\cel{4}r+}
}
